\documentclass[11pt]{article}
\usepackage[margin=1in]{geometry}
\usepackage{amsmath,amssymb,amsfonts,amsthm,mathtools,bm}
\usepackage[T1]{fontenc}
\usepackage{microtype}
\newtheorem{theorem}{Theorem}
\newtheorem{lemma}{Lemma}
\newtheorem{cor}{Corollary}
\newtheorem{defn}{Definition}
\newtheorem{principle}{Principle}
\newtheorem{statement}{Statement}
\theoremstyle{remark}
\newtheorem{statementr}{Statement}

\title{On Recovering a Wave Function from Bohmian Velocities}
\author{James Taylor}
\date{Working draft derived from the 2003 thesis}
\begin{document}
\maketitle

\section{On bijective 
correspondences between Bohmian velocity fields
  and wave functions} \label{a:ubvel}

The PDE for the Bohmian trajectories in Chapter \ref{c:algo}
 simplified greatly in
the situations in which the second derivative of $Q$ vanished.
 Indeed, we find in that situation that we can have
multiple initial wave functions which have the same Bohmian velocity
fields. They must have the same phase and their amplitudes must
satisfy $\frac{\ensuremath{\Delta} R_1}{R_1} = \frac{\ensuremath{\Delta} R_2}{R_2} +C$, where $C$ is a
 constant. It is possible to have multiple \ensuremath{L^{2}}\ solutions of this
 equation. These
include all  known cases in which the system has the
 property that the same velocity
field corresponds to different initial wave functions. 

If the second derivative of $Q$ does not vanish, then the PDE becomes
pretty intractable. The right hand side has to be the same in order to
produce the same quantum forces. But now we see that the requirement on
$R$ is more intensive
than just that equation above, which still must be satisfied. 
We therefore could speculate
that it can no longer happen. Although we do not see how to derive it
from this PDE, we do know of some other arguments. Two arguments have
been done by others; see \cite{S95, Bac03}. We shall give yet another
argument. Although lacking in rigor, it may give an appropriate
picture as to when and why the claim holds.


First, for two positive functions to satisfy the same stationary state
equation, we essentially need one of them to have codimension one with
respect to the other. This follows from energy arguments, i.e. the
ground state is positive and the only such wave function. If there were two,
then they could be combined in such a way as to lower the
energy. Actually, one first argues that they must have the same
energy; that is to say, there can only be one energy value with
positive solutions. This is easy; indeed, the functions must be
orthogonal if they are eigenstates for the same energy value. But two
positive functions cannot be orthogonal. That is essentially it. Then,
one can demonstrate that at the same eigenvalue, one can create states
of lesser energy. Therefore, if the potential, the one defining the stationary
state equation, is such that it has a
lowest energy state, then we are done.

Second, typically a wave function's nodal set has codimension 2 since it
requires the vanishing of both the real and imaginary
parts. Therefore, for a typical evolution, one would expect the Schr\"odinger\
evolution to instantly take a codimension 1 wave function\ into a codimension 2
wave function. There can, of course, be exceptions, but these must be
explained. We therefore expect that the full-time Bohmian velocity field is
generically in bijective correspondence with the initial wave function.

\clearpage
\bibliographystyle{abbrv}
\bibliography{references}
\end{document}
